\documentclass[11pt]{article}

\usepackage[a4paper,margin=1in]{geometry}
\usepackage{amsmath,amssymb,mathtools,bm}
\usepackage{enumitem}
\usepackage{booktabs}
\usepackage{hyperref}
\usepackage{xcolor}

\newcommand{\R}{\mathbb{R}}
\newcommand{\E}{\mathbb{E}}
\newcommand{\1}{\mathbf{1}}
\newcommand{\dd}{\mathrm{d}}
\newcommand{\diag}{\operatorname{diag}}
\newcommand{\atanh}{\operatorname{atanh}}
\newcommand{\atanTwo}{\operatorname{atan2}}
\newcommand{\sech}{\operatorname{sech}}
\newcommand{\MLE}{\mathrm{MLE}}
\newcommand{\Jac}{\mathcal{J}}

\title{Constrained Gibbs Sampling for Bayesian Inference\\from a Released Location--Scale MLE}
\author{}
\date{}

\begin{document}
\maketitle

\section{Purpose and scope}

This document develops a constrained Gibbs methodology for a one-dimensional location--scale family when the full data are latent and only the joint maximum likelihood estimator (MLE)
\[
    (\widehat\mu,\widehat\sigma)=(\mu^\star,\sigma^\star)
\]
is observed.  The main goals are:
\begin{enumerate}[label=(\roman*)]
    \item to identify the exact augmented target distribution;
    \item to explain why a local continuous update naturally uses \emph{three} observations rather than two;
    \item to reduce the three-observation update to a one-dimensional sampling problem;
    \item to exploit additional algebraic structure for the logistic and Student-$t_k$ families;
    \item to make explicit the distinction between conditioning on the score equations and conditioning on the actual MLE.
\end{enumerate}

The last point is essential.  A score root need not be the selected/global MLE, and even when the score-zero set and the MLE level set coincide as sets, the induced conditional measures generally differ by a Jacobian factor.  The formulas below therefore keep the MLE-conditioning correction explicit.

\section{Generic location--scale model}

Let
\[
    X_i\mid \theta \overset{\mathrm{iid}}{\sim}
    p_\theta(x)
    =\frac{1}{\sigma}
    f\!\left(\frac{x-\mu}{\sigma}\right),
    \qquad
    \theta=(\mu,\eta),\qquad \eta=\log\sigma,
\]
with prior density $\pi_0(\theta)$.  Write
\[
    r_i(\theta)=\frac{x_i-\mu}{\sigma},
    \qquad
    s(r)=\frac{\dd}{\dd r}\log f(r).
\]
For one observation,
\[
    \frac{\partial}{\partial \mu}\log p_\theta(x)
    =-\frac{1}{\sigma}s(r),
    \qquad
    \frac{\partial}{\partial \eta}\log p_\theta(x)
    =-\{1+r s(r)\}.
\]
Thus the joint MLE is characterized by the two estimating equations
\begin{align}
    C_1(x;\theta)
    &=\sum_{i=1}^n s(r_i(\theta))=0, \\
    C_2(x;\theta)
    &=\sum_{i=1}^n \{1+r_i(\theta)s(r_i(\theta))\}=0.
\end{align}
Define
\[
    C(x;\theta)
    =
    \begin{pmatrix}
    C_1(x;\theta)\\
    C_2(x;\theta)
    \end{pmatrix}.
\]
The observed release is
\[
    T(x)=\widehat\theta(x)=\theta^\star
    =(\mu^\star,\eta^\star),
    \qquad \eta^\star=\log\sigma^\star.
\]

\section{Score conditioning versus exact MLE conditioning}

\subsection{The score manifold}

The score-zero manifold at the released value is
\[
    \mathcal M_{\mathrm{score}}(\theta^\star)
    =\{x\in\R^n:C(x;\theta^\star)=0\}.
\]
It has codimension two under the usual full-rank regularity condition.

If $\theta^\star$ is an interior MLE, then necessarily
\[
    T(x)=\theta^\star
    \quad\Longrightarrow\quad
    C(x;\theta^\star)=0.
\]
The converse requires an additional MLE-selection condition: the score root must be the selected/global maximizer.  We therefore define
\[
    I_{\MLE}(x;\theta^\star)
    =\1\{\theta^\star\text{ is the selected MLE for }x\}.
\]
If uniqueness/globality has been established throughout the region of interest, this indicator is identically one and may be dropped.

\subsection{The MLE-conditioning Jacobian}

Assume locally that the estimator $T(x)$ is implicitly defined by
\[
    C(x;T(x))=0
\]
and that
\[
    H_C(x;\theta^\star)
    :=D_\theta C(x;\theta)\big|_{\theta=\theta^\star}
\]
is nonsingular.  Implicit differentiation gives
\[
    D_xT(x)
    =-H_C(x;\theta^\star)^{-1}
      D_xC(x;\theta^\star).
\]
Consequently, the coarea Jacobians satisfy
\[
    J_T(x)
    =\frac{J_C(x)}{|\det H_C(x;\theta^\star)|},
\]
where
\[
    J_C(x)
    =\sqrt{\det\!\left(D_xC\,D_xC^\mathsf T\right)}.
\]
Therefore, with respect to Hausdorff surface measure on the level set,
\begin{equation}
\boxed{
    p(x\mid \theta,\,T(x)=\theta^\star)
    \propto
    p_\theta(x)
    \frac{|\det H_C(x;\theta^\star)|}{J_C(x)}
    I_{\MLE}(x;\theta^\star).
}
\label{eq:exact-mle-manifold-target}
\end{equation}
By contrast, conditioning only on $C(x;\theta^\star)=0$ would omit the factor
$|\det H_C|$ and the MLE-selection indicator.

\paragraph{Important invariance.}
The estimating equations may be multiplied by any nonsingular matrix depending only on $\theta$: replacing $C$ by $A(\theta)C$ changes numerator and denominator in the same way at a root.  Hence convenient rescalings or sign changes of the two score equations do not change the conditional law.

\subsection{Explicit form of the $2\times2$ correction}

Let $r_i=(x_i-\mu)/\sigma$.  Since
\[
    C_1=\sum_i s(r_i),
    \qquad
    C_2=\sum_i \{1+r_i s(r_i)\},
\]
we have
\begin{equation}
\boxed{
H_C(x;\theta)
=
\begin{pmatrix}
-\sigma^{-1}\sum_i s'(r_i)
& -\sum_i r_i s'(r_i)\\[4pt]
-\sigma^{-1}\sum_i \{s(r_i)+r_i s'(r_i)\}
& -\sum_i r_i\{s(r_i)+r_i s'(r_i)\}
\end{pmatrix}.
}
\label{eq:generic-H}
\end{equation}
This matrix is additive over observations, so its value can be updated in $O(1)$ work when only three coordinates change.

\section{Augmented Gibbs sampler}

The augmented state is
\[
    (\theta,x),\qquad T(x)=\theta^\star.
\]
The target posterior is
\[
    p(\theta,x\mid T(x)=\theta^\star).
\]
A Gibbs iteration contains two blocks.

\subsection{Parameter block}

Conditional on the latent data $x$, the MLE event is deterministic and contributes no $\theta$-dependent factor.  Therefore
\begin{equation}
\boxed{
    p(\theta\mid x,T(x)=\theta^\star)
    \propto
    \pi_0(\theta)
    \prod_{i=1}^n
    \frac{1}{\sigma}
    f\!\left(\frac{x_i-\mu}{\sigma}\right).
}
\end{equation}
This is an ordinary two-dimensional Bayesian update and may be implemented with Metropolis--Hastings, slice sampling, HMC, or another valid kernel.

\subsection{Why the latent block should update triples}

At fixed $\theta^\star$, each observation contributes two numbers to the score constraint.  Define
\[
    y_i=\frac{x_i-\mu^\star}{\sigma^\star}
\]
and
\begin{equation}
\boxed{
    G(y)
    =
    \begin{pmatrix}
        s(y)\\
        1+y s(y)
    \end{pmatrix}.
}
\end{equation}
The released-MLE score constraint is
\[
    \sum_{i=1}^n G(y_i)=0.
\]

A pair update has two unknown coordinates and two constraints, so generically it has zero continuous degrees of freedom.  A triple update has
\[
    3\text{ variables}-2\text{ constraints}=1
\]
continuous degree of freedom.  Hence three observations are the natural analogue of the pair update used in a one-parameter location family.

Choose indices $i,j,\ell$.  Their current contribution is
\begin{equation}
\boxed{
    \Delta
    =G(y_i)+G(y_j)+G(y_\ell)\in\R^2.
}
\end{equation}
All other coordinates are held fixed, so the new triple must satisfy
\begin{equation}
\boxed{
    G(y_i')+G(y_j')+G(y_\ell')=\Delta.
}
\label{eq:triple-constraint}
\end{equation}
Preserving $\Delta$ preserves both global score equations exactly up to numerical tolerance.

\section{Exact one-dimensional conditional for a generic triple}
\label{sec:generic-sampling}

This section is the central sampling calculation.

\subsection{Density of one released-standardized coordinate}

During the latent update, the current model parameter is
\[
    \theta=(\mu,\sigma),
\]
whereas $(\mu^\star,\sigma^\star)$ is fixed by the released statistic.  Since
\[
    x=\mu^\star+\sigma^\star y,
\]
the density of $y$ under the current model parameter is
\begin{equation}
\boxed{
    \rho_\theta(y)
    =\frac{\sigma^\star}{\sigma}
      f\!\left(
      \frac{\mu^\star+\sigma^\star y-\mu}{\sigma}
      \right).
}
\label{eq:rho-y}
\end{equation}
The prefactor $\sigma^\star/\sigma$ is constant in $y$ and can be omitted in unnormalized calculations.

\subsection{Optional one-to-one scalar transformation}

Let
\[
    w=h_0(y),\qquad y=\chi(w)
\]
be a one-to-one transformation chosen for convenience.  Define
\[
    q_\theta(w)
    =\rho_\theta(\chi(w))\,|\chi'(w)|,
\]
and
\[
    \bar G(w)=G(\chi(w)).
\]
We now use $w_i$ as the free coordinate.

\subsection{Fix one coordinate and reconstruct the other two}

Set
\[
    t=w_i'.
\]
The remaining pair must satisfy
\begin{equation}
    \bar G(u)+\bar G(v)
    =\Delta-\bar G(t).
\label{eq:pair-reconstruction-generic}
\end{equation}
For a given $t$, let
\[
    \mathcal S(t)
    =\{(u_b(t),v_b(t)):b\in\mathcal B(t)\}
\]
denote all isolated labeled solutions of \eqref{eq:pair-reconstruction-generic}.

For each branch $b$, define the pair Jacobian
\begin{equation}
\boxed{
    J_b(t)
    =
    \left|
    \det
    \begin{pmatrix}
        \bar G_1'(u_b) & \bar G_1'(v_b)\\
        \bar G_2'(u_b) & \bar G_2'(v_b)
    \end{pmatrix}
    \right|.
}
\label{eq:pair-jac}
\end{equation}
This is the ordinary change-of-variables Jacobian for the map
\[
    (u,v)\mapsto \bar G(u)+\bar G(v)\in\R^2.
\]

\subsection{The scalar density that must actually be sampled}

Let $x_b(t)$ denote the full latent dataset obtained by inserting the candidate triple associated with branch $b$.  Then, with respect to $\dd t$, the exact MLE-conditioned scalar target is
\begin{equation}
\boxed{
\pi(t\mid x_{-(i,j,\ell)},\theta,T=\theta^\star)
\propto
q_\theta(t)
\sum_{b\in\mathcal B(t)}
\frac{
q_\theta(u_b(t))q_\theta(v_b(t))
\,|\det H_C(x_b(t);\theta^\star)|
\,I_{\MLE}(x_b(t);\theta^\star)
}{J_b(t)}.
}
\label{eq:master-scalar-target}
\end{equation}

This equation is the main correctness criterion for the local update.

After sampling $t$, the solution branch is sampled conditionally with probability
\begin{equation}
\boxed{
\Pr(B=b\mid t,\cdots)
\propto
\frac{
q_\theta(u_b)q_\theta(v_b)
\,|\det H_C(x_b;\theta^\star)|
\,I_{\MLE}(x_b;\theta^\star)
}{J_b}.
}
\label{eq:branch-weight-generic}
\end{equation}
Finally transform $(t,u_b,v_b)$ back to $(y_i',y_j',y_\ell')$ and then to $x$.

\paragraph{Interpretation.}
There are four conceptually distinct ingredients in \eqref{eq:master-scalar-target}:
\begin{enumerate}[label=(\alph*)]
    \item the model density of the three latent observations;
    \item the Jacobian induced by any scalar reparameterization;
    \item the inverse pair-map Jacobian $1/J_b$, which arises because two variables are being conditioned on two equations;
    \item the MLE-conditioning correction $|\det H_C|$ and, when necessary, the global-MLE selection indicator.
\end{enumerate}
Omitting any nonconstant factor changes the target distribution.

\paragraph{How to sample $t$.}
Equation \eqref{eq:master-scalar-target} is one-dimensional.  One may use an exact rejection/inverse-CDF draw when practical, or a valid one-dimensional MCMC kernel such as slice sampling or Metropolis--Hastings.  In the latter case the overall method is technically Metropolis-within-Gibbs, but it still targets the correct augmented posterior if the scalar kernel leaves \eqref{eq:master-scalar-target} invariant.

\section{Logistic location--scale family}

\subsection{Score simplification}

For the standard logistic density
\[
    f(y)=\frac{1}{4\cosh^2(y/2)},
\]
we have
\[
    s(y)=\frac{\dd}{\dd y}\log f(y)
    =-\tanh(y/2).
\]
Define the one-to-one transformation
\begin{equation}
\boxed{
    z=\tanh(y/2),
    \qquad
    y=2\atanh(z),
    \qquad z\in(-1,1).
}
\end{equation}
Then the two score equations are equivalently
\begin{align}
    \sum_i z_i &=0,\\
    \sum_i h(z_i)&=0,
\end{align}
where
\begin{equation}
\boxed{
    h(z)=2z\atanh(z)-1.
}
\end{equation}
Thus we may use
\[
    \bar G(z)
    =
    \begin{pmatrix}
        z\\h(z)
    \end{pmatrix}.
\]
The transformed single-coordinate density is
\begin{equation}
\boxed{
    q_\theta(z)
    =\rho_\theta(2\atanh z)\frac{2}{1-z^2},
    \qquad z\in(-1,1).
}
\end{equation}

\subsection{Triple reconstruction reduces to one scalar root}

For a selected triple, store
\[
    \Delta
    =\bar G(z_i)+\bar G(z_j)+\bar G(z_\ell).
\]
Choose the free value
\[
    t=z_i'.
\]
Write
\[
    A(t)=\Delta_1-t,
    \qquad
    B(t)=\Delta_2-h(t).
\]
The remaining pair must satisfy
\[
    z_j'+z_\ell'=A(t),
    \qquad
    h(z_j')+h(z_\ell')=B(t).
\]
Set
\[
    z_\ell'=A(t)-z_j'.
\]
Then $z_j'$ solves the one-dimensional equation
\begin{equation}
\boxed{
    F_t(z)
    :=h(z)+h(A(t)-z)-B(t)=0,
}
\label{eq:logistic-root}
\end{equation}
with feasible interval
\[
    z\in
    \bigl(\max\{-1,A(t)-1\},\,\min\{1,A(t)+1\}\bigr).
\]
Because
\[
    h''(z)=\frac{4}{(1-z^2)^2}>0,
\]
$F_t$ is strictly convex and symmetric about $A(t)/2$.  Hence there are at most two roots, which are the two label-swapped solutions $(z_j',z_\ell')$ and $(z_\ell',z_j')$.  In practice one root may be found on one side of $A(t)/2$ with a robust one-dimensional solver such as Brent's method; the other is obtained by swapping.

\subsection{Exact scalar target for logistic}

Since
\[
    \bar G'(z)
    =
    \begin{pmatrix}
        1\\h'(z)
    \end{pmatrix},
\]
the pair Jacobian is
\begin{equation}
\boxed{
    J_{\mathrm{log}}(z_j,z_\ell)
    =|h'(z_\ell)-h'(z_j)|,
}
\end{equation}
where
\[
    h'(z)
    =2\atanh(z)+\frac{2z}{1-z^2}.
\]
The two label-swapped branches have the same weight, so the scalar density simplifies, up to an irrelevant factor of two, to
\begin{equation}
\boxed{
\pi_{\mathrm{log}}(t)
\propto
\frac{
q_\theta(t)
q_\theta(z_j(t))
q_\theta(z_\ell(t))
}{|h'(z_\ell(t))-h'(z_j(t))|}
\,|\det H_C(x(t);\theta^\star)|
\,I_{\MLE}(x(t);\theta^\star).
}
\label{eq:logistic-scalar-target}
\end{equation}
After drawing $t$, the two labeled solutions may be selected with probability $1/2$ each.

For the MLE correction, one may evaluate \eqref{eq:generic-H} using
\[
    s(y)=-\tanh(y/2),
    \qquad
    s'(y)=-\frac12\sech^2(y/2).
\]

\paragraph{Computational summary for logistic.}
Each candidate evaluation of the scalar target requires one one-dimensional pair reconstruction through \eqref{eq:logistic-root}, followed by evaluation of the three transformed densities, the pair Jacobian, and the $2\times2$ MLE-conditioning determinant.

\section{Student-$t_k$ location--scale family}

\subsection{Joint score equations}

For the standard Student-$t_k$ density,
\[
    f_k(y)
    \propto
    \left(1+\frac{y^2}{k}\right)^{-(k+1)/2},
\]
we have
\[
    s(y)=-(k+1)\frac{y}{k+y^2}.
\]
The location and scale score equations are equivalently
\begin{align}
    \sum_{i=1}^n \frac{y_i}{k+y_i^2}&=0,
    \label{eq:student-loc-score}\\
    \sum_{i=1}^n \frac{y_i^2}{k+y_i^2}&=\frac{n}{k+1}.
    \label{eq:student-scale-score}
\end{align}

\subsection{Angle transformation: the two constraints become a vector sum}

Define
\begin{equation}
\boxed{
    \phi
    =2\arctan\!\left(\frac{y}{\sqrt{k}}\right),
    \qquad
    y=\sqrt{k}\tan(\phi/2),
    \qquad
    \phi\in(-\pi,\pi).
}
\end{equation}
Then
\[
    \sin\phi
    =\frac{2\sqrt{k}\,y}{k+y^2},
    \qquad
    \cos\phi
    =\frac{k-y^2}{k+y^2}.
\]
Therefore \eqref{eq:student-loc-score}--\eqref{eq:student-scale-score} become
\begin{equation}
\boxed{
    \sum_{i=1}^n \sin\phi_i=0,
    \qquad
    \sum_{i=1}^n \cos\phi_i
    =n\frac{k-1}{k+1}.
}
\end{equation}
Define the unit vector
\begin{equation}
\boxed{
    u(\phi)
    =
    \begin{pmatrix}
        \cos\phi\\
        \sin\phi
    \end{pmatrix}.
}
\end{equation}
Then the entire released-MLE score constraint is simply
\begin{equation}
\boxed{
    \sum_{i=1}^n u(\phi_i)
    =
    \begin{pmatrix}
    n(k-1)/(k+1)\\0
    \end{pmatrix}.
}
\label{eq:student-resultant}
\end{equation}
Thus the Student location--scale manifold becomes a fixed-resultant problem for unit vectors on the circle.

\subsection{Triple reconstruction is closed form}

For a selected triple, preserve its current resultant
\begin{equation}
\boxed{
    R
    =u(\phi_i)+u(\phi_j)+u(\phi_\ell)\in\R^2.
}
\end{equation}
Choose the free angle
\[
    t=\phi_i'.
\]
The remaining two vectors must sum to
\[
    v(t)=R-u(t).
\]
Let
\[
    r(t)=\|v(t)\|.
\]
A pair of unit vectors with resultant $v(t)$ exists if and only if
\[
    0\le r(t)\le 2.
\]
For the regular case $0<r(t)<2$, define
\begin{equation}
\boxed{
    \beta(t)=\atanTwo(v_2(t),v_1(t)),
    \qquad
    \gamma(t)=\arccos\!\left(\frac{r(t)}{2}\right).
}
\end{equation}
Then the two labeled solutions are
\begin{equation}
\boxed{
    \phi_j'=\beta+\gamma,
    \qquad
    \phi_\ell'=\beta-\gamma,
}
\label{eq:student-angle-pair}
\end{equation}
and the label-swapped version.  Angles are wrapped back to $(-\pi,\pi)$.

Hence, unlike the generic or logistic construction, \emph{no root solve is required to reconstruct the remaining Student pair}.  The only continuously sampled quantity is the single free angle $t$.

\subsection{Density in angle space}

From
\[
    y=\sqrt{k}\tan(\phi/2),
\]
we obtain
\[
    \left|\frac{\dd y}{\dd\phi}\right|
    =\frac{\sqrt{k}}{2}\sec^2(\phi/2).
\]
Therefore the density of one angle under the current Gibbs parameter $\theta=(\mu,\sigma)$ is
\begin{equation}
\boxed{
    q_\theta(\phi)
    =\rho_\theta\!\left(\sqrt{k}\tan(\phi/2)\right)
      \frac{\sqrt{k}}{2}\sec^2(\phi/2).
}
\label{eq:student-angle-density}
\end{equation}

\subsection{Exact scalar target for Student-$t_k$}

For the pair map
\[
    (\phi_j,\phi_\ell)
    \mapsto
    u(\phi_j)+u(\phi_\ell),
\]
the absolute Jacobian determinant is
\begin{equation}
\boxed{
    J_{\mathrm{Stud}}
    =|\sin(\phi_j-\phi_\ell)|
    =|\sin(2\gamma)|.
}
\end{equation}
The label-swapped solutions again have identical weights.  Thus, up to an irrelevant factor of two, the exact MLE-conditioned scalar density is
\begin{equation}
\boxed{
\pi_{\mathrm{Stud}}(t)
\propto
\frac{
q_\theta(t)
q_\theta(\beta(t)+\gamma(t))
q_\theta(\beta(t)-\gamma(t))
}{|\sin(2\gamma(t))|}
\,|\det H_C(x(t);\theta^\star)|
\,I_{\MLE}(x(t);\theta^\star),
}
\label{eq:student-scalar-target}
\end{equation}
for feasible $t$ satisfying $0<\|R-u(t)\|<2$.

After drawing $t$:
\begin{enumerate}
    \item compute $\beta(t)$ and $\gamma(t)$;
    \item reconstruct the other two angles by \eqref{eq:student-angle-pair};
    \item randomly assign the two reconstructed angles to labels $j$ and $\ell$ with probability $1/2$ each;
    \item transform back through
    \[
        y_r'=\sqrt{k}\tan(\phi_r'/2),
        \qquad
        x_r'=\mu^\star+\sigma^\star y_r'.
    \]
\end{enumerate}
The local resultant $R$ is then preserved exactly up to floating-point error, so both score constraints remain satisfied.

For the MLE correction, evaluate \eqref{eq:generic-H} with
\[
    s(y)=-(k+1)\frac{y}{k+y^2},
\]
and
\begin{equation}
\boxed{
    s'(y)
    =-(k+1)\frac{k-y^2}{(k+y^2)^2}.
}
\end{equation}

\paragraph{Singular geometric cases.}
The cases $r(t)=0$ and $r(t)=2$ are singular for the pair reconstruction: at $r=0$ the remaining two unit vectors can be any opposite pair, and at $r=2$ they coincide.  These are lower-dimensional boundary cases and have probability zero under a regular continuous target, but numerical implementations should guard against evaluating exactly at them.

\section{What is actually being sampled?}

For both logistic and Student-$t_k$, the essential latent move is:
\begin{quote}
\emph{Choose a random triple, preserve its two-dimensional score contribution, and sample one scalar free coordinate from the exact one-dimensional conditional density induced by the model likelihood, the pair reconstruction Jacobian, and the MLE-conditioning factor.  The remaining two coordinates are then reconstructed from the two preserved constraints.}
\end{quote}

For logistic, the remaining pair is obtained from one scalar root solve.  For Student-$t_k$, the angle transformation makes the remaining pair available in closed form.

The exactness test is therefore not merely whether the new triple satisfies the two score equations.  The update must also use the correct scalar density \eqref{eq:master-scalar-target}, including its Jacobian and MLE-conditioning terms.

\section{Alternative affine-normalization proposal}

Location--scale MLEs are affine equivariant.  For any unconstrained dataset $u$, define
\[
    a=\widehat\mu(u),
    \qquad
    b=\widehat\sigma(u),
\]
and standardize
\[
    z_i=\frac{u_i-a}{b}.
\]
Then
\[
    \widehat\mu(z)=0,
    \qquad
    \widehat\sigma(z)=1.
\]
Hence
\begin{equation}
\boxed{
    x_i'
    =\mu^\star+\sigma^\star
      \frac{u_i-\widehat\mu(u)}{\widehat\sigma(u)}
}
\end{equation}
automatically satisfies
\[
    (\widehat\mu(x'),\widehat\sigma(x'))
    =(\mu^\star,\sigma^\star).
\]

This is a useful way to generate \emph{constraint-satisfying proposals}, but it is not by itself an exact Gibbs update.  The map from $u$ to its standardized shape is many-to-one: all affine datasets
\[
    a\mathbf 1+bz,
    \qquad a\in\R,\quad b>0,
\]
produce the same standardized configuration.  Therefore the induced proposal density on the MLE manifold contains an affine-orbit Jacobian/integration factor.  Affine normalization is best viewed as a possible global Metropolis--Hastings proposal unless that induced density is derived explicitly.

\section{Algorithmic template}

\begin{enumerate}[label=\textbf{Step \arabic*:},leftmargin=*]
    \item \textbf{Initialize.} Choose $(\theta^{(0)},x^{(0)})$ with
    \[
        T(x^{(0)})=\theta^\star.
    \]

    \item \textbf{Parameter update.} Sample/update
    \[
        \theta^{(t)}\sim p(\theta\mid x^{(t-1)})
        \propto \pi_0(\theta)\prod_i p_\theta(x_i^{(t-1)}).
    \]

    \item \textbf{Randomly partition or repeatedly select triples.}

    \item \textbf{For each selected triple $(i,j,\ell)$:}
    \begin{enumerate}[label=(\alph*)]
        \item compute its fixed two-dimensional score contribution $\Delta$;
        \item choose the family-specific free coordinate $t$;
        \item enumerate/reconstruct every regular pair solution for the remaining two coordinates;
        \item evaluate the exact scalar target, including model densities, transformation Jacobians, inverse pair Jacobian, $|\det H_C|$, and the MLE-selection condition when needed;
        \item sample/update $t$ with a valid one-dimensional kernel;
        \item sample the reconstruction branch using its conditional weight;
        \item write the new triple back to $x$.
    \end{enumerate}

    \item \textbf{Diagnostics.} Verify after every sweep that
    \[
        \|C(x;\theta^\star)\|
    \]
    is at numerical zero and, when exact MLE conditioning is claimed, periodically verify that $\theta^\star$ is indeed the selected MLE of sampled latent datasets.
\end{enumerate}

\section{Implementation correctness checklist}

A coding implementation should not be declared correct until the following are explicitly verified:
\begin{enumerate}[leftmargin=*]
    \item Both location and scale score constraints are preserved after every latent update.
    \item Three coordinates, rather than two, are used for a one-dimensional continuous local move.
    \item The scalar free coordinate is sampled from the correct one-dimensional conditional, not uniformly on the feasible set.
    \item Every regular reconstruction branch is included with the correct inverse-Jacobian weight.
    \item Any change of variables includes its ordinary Jacobian.
    \item Exact MLE conditioning includes $|\det H_C(x;\theta^\star)|$ relative to score conditioning.
    \item If score equations can have non-MLE roots, candidates are restricted to the selected/global-MLE branch or the regime is explicitly treated as unresolved.
    \item For logistic, the pair root solve accounts for the full feasible interval and both label-swapped solutions.
    \item For Student-$t_k$, the angle reconstruction uses only candidates with $0<\|R-u(t)\|<2$ and handles periodic angle wrapping correctly.
    \item All densities and Jacobian factors are evaluated in the log domain near singular or heavy-tail regions.
\end{enumerate}

\section{Summary of the simplifications}

\begin{center}
\begin{tabular}{@{}llll@{}}
\toprule
Family & Free coordinate & Remaining pair & Main pair Jacobian \\
\midrule
Generic location--scale & one scalar & solve two equations & $|\det[\bar G'(u),\bar G'(v)]|$ \\
Logistic & $z_i\in(-1,1)$ & one scalar root & $|h'(z_\ell)-h'(z_j)|$ \\
Student-$t_k$ & $\phi_i\in(-\pi,\pi)$ & closed-form angles & $|\sin(\phi_j-\phi_\ell)|$ \\
\bottomrule
\end{tabular}
\end{center}

The core generalization is therefore
\[
\boxed{
\text{one score constraint + pair update}
\quad\longrightarrow\quad
\text{two score constraints + triple update}.
}
\]
The logistic transformation reduces pair reconstruction to a one-dimensional root problem, while the Student-$t_k$ angle transformation converts the two score equations into a fixed vector-resultant constraint on the unit circle and removes the reconstruction root solve entirely.

\section*{References}

\begin{thebibliography}{9}

\bibitem{graham2022}
M. M. Graham, A. H. Thiery, and A. Beskos.
\newblock Manifold Markov chain Monte Carlo methods for Bayesian inference in diffusion models.
\newblock \emph{Journal of the Royal Statistical Society: Series B}, 84:1229--1256, 2022.

\bibitem{byrne2013}
S. Byrne and M. Girolami.
\newblock Geodesic Monte Carlo on embedded manifolds.
\newblock \emph{Scandinavian Journal of Statistics}, 40:825--845, 2013.

\bibitem{lelievre2019}
T. Leli\`evre, M. Rousset, and G. Stoltz.
\newblock Hybrid Monte Carlo methods for sampling probability measures on submanifolds.
\newblock \emph{Numerische Mathematik}, 143:379--421, 2019.

\end{thebibliography}

\end{document}
